
%----------------------------------------
\begin{frame}<0>[fragile]
  \frametitle{A mini SMT lecture} \small

\begin{block}{SMT problem: constraint on an integer variable}
\[
\begin{aligned}
& x : \intg \\
& x > 10
\end{aligned}
\]
\end{block}

\pause

\[
[x \mapsto 11]
\]

\pause

\begin{block}{SMT problem: constraint on a function}
\[
\begin{aligned}
& f : \mathbb{Z} \to \mathbb{Z} \\
& \forall x, y \in \mathbb{Z},\ f(x) = f(y) \Rightarrow x = y
\end{aligned}
\]
\end{block}

\pause

\[
[f \mapsto \lambda x. x]
\]

\end{frame}


%----------------------------------------
\begin{frame}<0>[fragile]
  \frametitle{Constraint on a function in SMT-LIB} \small

\begin{block}{SMT problem: constraint on a function}
\begin{verbatim}
    (declare-fun f (Int) Int)
    (assert
      (forall ((x Int) (y Int))
        (=> (= (f x) (f y)) (= x y))
      )
    )
\end{verbatim}
\end{block}

\begin{verbatim}
    sat
    (model
      (define-fun f ((x Int)) Int
        x)
    )
\end{verbatim}

\end{frame}


%----------------------------------------
\begin{frame}<0>[fragile]
  \frametitle{Representing a test as a function from time to commands} \small

{\bf The SMT type of tests:}

\begin{lstlisting}[language=python,frame=single]
test: Function = Function('test', IntSort(), Command)
\end{lstlisting}

\pause

{\bf The SMT type of commands:}

{\scriptsize
\begin{verbatim}
(declare-datatypes ()
  (
    (Command
      (ROTATE (ROTATE_number Int)(ROTATE_time Int)(ROTATE_angle Real))
      (GOTO (GOTO_number Int)(GOTO_time Int)(GOTO_x Real)(GOTO_y Real))
      (MOVE (MOVE_number Int)(MOVE_time Int)
            (MOVE_dir direction)(MOVE_distance Real))
      (PIC (PIC_number Int)(PIC_time Int)
           (PIC_quality image_quality)(PIC_images Int))
      (STORE (STORE_number Int)(STORE_time Int)
             (STORE_file String)(STORE_images Int))
      (SEND (SEND_number Int)(SEND_time Int)(SEND_file String))
      (COLLECT (COLLECT_number Int)(COLLECT_time Int)
               (COLLECT_file String)(COLLECT_sample String))
      (SCRIPT (SCRIPT_number Int)(SCRIPT_time Int)
              (SCRIPT_script String)(SCRIPT_file String))
    )
  )
)    
\end{verbatim}
}

\end{frame}


%----------------------------------------
\begin{frame}<0>[fragile]
  \frametitle{Modeling temporal logic in first-order logic on a finite trace} \small

{\bf Consider the \underline{propositional} temporal formula:} $\Box (p \rightarrow \Diamond q)$

\pause

\vspace{0.2cm}
{\bf One can model it in first-order logic with quantifiers:}

\[
\forall i \in \{0..k\!-\!1\},\ p(test(i)) 
  \rightarrow 
    \exists j \in \{i..k\!-\!1\}, q(test(j)))
\]

\pause

\vspace{0.2cm}
{\bf For efficiency we instead ``unfold'' the quantifiers (here for trace length 3):} 

\begin{align*}
&p(test(0)) \rightarrow (q(test(0)) \vee q(test(1)) \vee q(test(2))) \\
\wedge\quad &p(test(1)) \rightarrow (q(test(1)) \vee q(test(2))) \\
\wedge\quad &p(test(2)) \rightarrow q(test(2))
\end{align*}

\pause

\vspace{0.2cm}
{\bf And a rewrite, recalling that $p \rightarrow q \equiv \neg p \vee q$:}

\begin{align*}
&\neg p(test(0)) \vee q(test(0)) \vee q(test(1)) \vee q(test(2)) \\
\wedge\quad &\neg p(test(1)) \vee q(test(1)) \vee q(test(2)) \\
\wedge\quad &\neg p(test(2)) \vee q(test(2))
\end{align*}

\end{frame}


%----------------------------------------
\newcommand{\isM}[1]{\mathtt{isM}(#1)}
\newcommand{\isR}[1]{\mathtt{isR}(#1)}
\newcommand{\Mdir}[1]{\mathtt{Mdir}(#1)}
\newcommand{\Mdist}[1]{\mathtt{Mdist}(#1)}
\newcommand{\Rangle}[1]{\mathtt{Rangle}(#1)}

\begin{frame}<0>[fragile]
  \frametitle{Let's now look at the more complicated property p5} \small

\begin{lstlisting}[language=fuzzdsl,frame=single]
  rule p5: # a backward move leads to a future rotation
    always [MOVE(dir=direction.backwards,distance=d?)] 
      eventually <<ROTATE(angle=a?)>> 
                     ((d <= 1 and a = 45) or (d > 1 and a = 20))  
\end{lstlisting}

\pause

\vspace{-0.3cm}

First we define some shorthands:
\hfill
\begin{minipage}{0.55\linewidth}
{\scriptsize
\begin{align*}
\isM{i}    &= \mathsf{is}(\mathsf{MOVE}, \mathsf{test}(i)) \\
\isR{i}    &= \mathsf{is}(\mathsf{ROTATE}, \mathsf{test}(i)) \\
\Mdir{i}   &= \mathsf{MOVE\_dir}(\mathsf{test}(i)) \\
\Mdist{i}  &= \mathsf{MOVE\_distance}(\mathsf{test}(i)) \\
\Rangle{i} &= \mathsf{ROTATE\_angle}(\mathsf{test}(i)) \\
\end{align*}
}
\end{minipage}

\pause

\vspace{-0.3cm}

\begin{block}{The SMT formula generated}
{\scriptsize
\[
\begin{aligned}
\bigwedge_{i=0}^{n-1} \Bigl(
   & \neg\bigl(\isM{i} \wedge \Mdir{i}=backward\bigr) \\[0.5ex]
   & \;\;\vee\;
     \bigvee_{j=i}^{n-1} \Bigl(
        \isR{j} \wedge
        \bigl(
           (\Mdist{i}\le 1 \wedge \Rangle{j}=45) 
           \vee\;
           (\Mdist{i} > 1 \wedge \Rangle{j}=20)
        \bigr)
     \Bigr)
\Bigr)
\end{aligned}
\]
}
\end{block}

\end{frame}

%----------------------------------------
\begin{frame}<0>[fragile]
  \frametitle{Show me the code} \small

\begin{block}{The until formula}
\[
\varphi_1 \until \varphi_2
\]
\end{block}

\begin{lstlisting}[language=python,frame=single]
@dataclass
class LTLUntil(LTLFormula):
  left: LTLFormula
  right: LTLFormula

  def to_smt(self, env: Environment, t: int, end_time: int) -> BoolRef: ...

  def evaluate(self, env: Environment, test: Test, index: int) -> bool: ...
\end{lstlisting}

\end{frame}


%----------------------------------------
\begin{frame}<0>[fragile]
  \frametitle{The to\_smt method} \small

\begin{block}{Follows this expansion approach}
\[
\varphi_1 \until \varphi_2 = 
  \varphi_2(0) 
  \vee (\varphi_2(1) \wedge \varphi_1(0))  
  \vee (\varphi_2(2) \wedge \varphi_1(0) \wedge \varphi_1(1))
  \vee ...
\]
\end{block}

\begin{lstlisting}[language=python,frame=single]
@dataclass
class LTLUntil(LTLFormula):
  left: LTLFormula
  right: LTLFormula

  def to_smt(self, env: Environment, t: int, end_time: int) -> BoolRef:
    return Or(
      [
        And(
          self.right.to_smt(env, t_prime, end_time),
          And([self.left.to_smt(env, t_i, end_time) 
                for t_i in range(t, t_prime)])
        )
        for t_prime in range(t, end_time)
      ])
\end{lstlisting}

\end{frame}


%----------------------------------------
\begin{frame}<0>[fragile]
  \frametitle{The evaluate method} \small

\begin{block}{Follows the semantics}
% \[
% (\gamma,\sigma,j) \models \varphi_2 \;\text{{\bf for some}}\; i \leq j < |\sigma|
% \;\text{{\bf and for all}}\; i \le k < j: \; (\gamma,\sigma,k) \models \varphi_1
% \]

\[
\varphi_1 \until \varphi_2 = \varphi_2 \vee (\varphi_1 \wedge \Next (\varphi_1 \until \varphi_2))
\]
\end{block}

\begin{lstlisting}[language=python,frame=single]
@dataclass
class LTLUntil(LTLFormula):
  left: LTLFormula
  right: LTLFormula

  def evaluate(self, env: Environment, test: Test, index: int) -> bool:
    if within(index, test):
      return self.right.evaluate(env, test, index) or (
        self.left.evaluate(env, test, index) 
          and 
        self.evaluate(env, test, index + 1)
      )
    return False
\end{lstlisting}

\pause

Since $\Diamond \varphi = \true \until \varphi\tt$ one can easily see that
the temporal formula $\Diamond \varphi$ will generate an SMT formula of the form 
$\varphi(0)  \vee \varphi(1) \vee \varphi(2)  \vee ... \vee \varphi(k-1)$.


\end{frame}
